% GFT: global algebraic tensor calculation, original full support carrier.
\section{Actual tensor powers over the entire Split-Zero coefficient semiring}
\label{sec:global-full-lattice-tensor}

\paragraph{Objects, sources, and scope.}
The coefficient object is the original full lattice construction in
\href{https://github.com/KokunoYumeto/zeta-function-research-reader/blob/a2851f9eb352e7e4376bc629de86fa4ed9c1c434/workbenches/splitzero-tandem/foundations/split-support-geometry-v11/split_support_geometry_arithmetic_curve_v11.tex}{\emph{Split Support Geometry}, v11, Definition \texttt{def:lattice-split}},
retained under its source byline The Clankers. The programme proof
FLH1--4, in \nolinkurl{FULL_LATTICE_NULL_HOMOLOGY.tex}, specifies its
null differential; those definitions are repeated below.
For the arithmetic application the human source is Alain Connes,
Caterina Consani and Matilde Marcolli,
\href{https://arxiv.org/abs/math/0703392v1}{\emph{The Weil proof and
the geometry of the adeles class space}}, author source
\nolinkurl{Yuri70v2.tex}, locators \texttt{rangecVdef},
\texttt{vanishV}, \texttt{fsharpinv}, \texttt{TraceformulaH1}
and \texttt{tracepairing}. SWQ1--24 constructs the actual closed
quotient from those definitions. GTK13,24,27--28 proves the
global mixed kernel and its original arithmetic trace, retaining
every prime, endpoint, real-place and conductor term.
The argument here proves the universal tensor map over the original
semiring, its entire quotient congruence and its full tensor pairing.
No completion of an algebraic tensor product, positive metric,
cohomological purity, or identification of a trace radical with
an adelic image is assumed.

\subsection{The universal tensor, including every zero label}

Let $L$ be any nontrivial bounded distributive lattice, and retain
\begin{equation}\tag{GFT1}
\begin{gathered}
 S=G_L(\mathbb C)=\{(0,\lambda):\lambda\in L\}
                       \cup(\mathbb C\times\{1_L\}),\\
 (a,\lambda)+(b,\mu)=(a+b,\lambda\vee\mu),\qquad
 (a,\lambda)(b,\mu)=(ab,\lambda\wedge\mu),\\
 z_\lambda=(0,\lambda),\quad \tau=z_{0_L},\quad
 e=z_{1_L},\quad \widehat a=(a,1_L),\quad \widehat0=e.
\end{gathered}
\end{equation}
For every complex vector space $V$ put
\begin{equation}\tag{GFT2}
 X(V)=G_L(V)=\{(0,\lambda):\lambda\in L\}
                          \cup(V\times\{1_L\}),\qquad
 (a,\eta)(v,\lambda)=(av,\eta\wedge\lambda).
\end{equation}
Addition uses the same join as GFT1. In particular $(0,1_L)$ is
not the additive identity of this semimodule. Every linear map
$A:V\to W$ has its actual lift
$A^\uparrow(v,\lambda)=(Av,\lambda)$; the lift of the zero
linear map is $e$ times the identity-label map, not the constant
map with value $\tau$.

The tensor product of $S$-semimodules below means the commutative
monoid on symbols $x\otimes y$ modulo additivity in both variables,
$\tau\otimes y=x\otimes\tau=0$, and
$(sx)\otimes y=x\otimes(sy)$, with its induced $S$ action.
This presentation proves its balanced biadditive universal property.
There is a canonical isomorphism
\begin{equation}\tag{GFT3}
 \Phi_{V,W}:X(V)\otimes_S X(W)\longrightarrow
                 X(V\otimes_{\mathbb C}W),\qquad
 (v,\lambda)\otimes(w,\mu)\longmapsto
                 (v\otimes w,\lambda\wedge\mu).
\end{equation}
This holds for the entire vector spaces, with no finite-dimensional
restriction. Nonzero output amplitude requires both nonzero input
amplitudes, hence both top labels. Distributivity of meet over join
proves biadditivity. The scalar law and associativity of meet prove
balancing. Thus the displayed formula defines an $S$-linear map.

Here is a proof that this map has the full universal property, so
that it has not silently imposed additional label identifications.
Let $P$ be any $S$-semimodule, and let
$b:X(V)\times X(W)\to P$ be balanced and biadditive. Write
$0_V^+=(0,1_L)$, $0_W^+=(0,1_L)$, and
\begin{equation}\tag{GFT4}
 t=b(0_V^+,0_W^+),\qquad P_t=\{p\in P:ep=t\}.
\end{equation}
One has $et=t$ and $t+t=t$. For any $p\in P_t$, the semiring
identities $1_S+e=1_S$ and $1_S+\widehat{-1}=e$ give
$p+t=p$ and $p+\widehat{-1}p=t$. The set $P_t$ is closed under
addition because $t+t=t$, and under every $\widehat a$ because
$\widehat a e=e$. These facts make $P_t$ a complex vector space,
with zero vector $t$ and scalar multiplication $a p=\widehat a p$.
The vector-space axioms are the semimodule axioms, including the
zero-scalar axiom $\widehat0 p=ep=t$.

Balancing by $e$ gives, for every $v,w$,
\begin{equation}\tag{GFT5}
 b(0_V^+,(w,1_L))=t=b((v,1_L),0_W^+),\qquad
 e b((v,1_L),(w,1_L))=t.
\end{equation}
Consequently $(v,w)\mapsto b((v,1_L),(w,1_L))$ is a
complex-bilinear map into $P_t$. The ordinary vector-space tensor
universal property gives a unique linear map
$u:V\otimes_{\mathbb C}W\to P_t$. Every lower-label value of
$b$ is forced by balancing:
\begin{equation}\tag{GFT6}
 b((0,\lambda),(w,\mu))=z_{\lambda\wedge\mu}t
 \quad\hbox{when }w=0\hbox{ or }\mu=1_L.
\end{equation}
These are all possible second inputs. For $\mu=1_L$ use
$(0,\lambda)=z_\lambda0_V^+$ and GFT5. For $w=0$ use
$(0,\mu)=z_\mu0_W^+$ and balance once more. The corresponding
identity holds when the second amplitude is zero.

Define the forced extension of $u$ by
\begin{equation}\tag{GFT7}
 U(z,1_L)=u(z),\qquad U(0,\lambda)=z_\lambda t.
\end{equation}
The formulas agree at $(0,1_L)$ since $u(0)=t=et$.
They are additive on lower labels because
$z_\lambda+z_\mu=z_{\lambda\vee\mu}$. On two top inputs
additivity is that of $u$. On a top and a lower input, if
$p=u(z)$ then
$p+z_\lambda t=(1_S+z_\lambda e)p=p$.
Scalar compatibility for a nonzero scalar is the linearity of
$u$ and $\widehat a z_\lambda=z_\lambda$. For a zero scalar
$z_\eta$ it is $z_\eta p=z_\eta ep=z_\eta t$, with the
displayed meet when the input label is lower. This proves every
semimodule law for $U$ and $U\Phi=b$. Uniqueness follows because
ordinary pure tensors span the top vector space and
$(0,\lambda)$ is the image of $0_V^+\otimes(0,\lambda)$.
Applying the two universal properties proves GFT3 is an isomorphism.

Iteration gives the exact $k$-fold isomorphism, for every $k\ge1$,
\begin{equation}\tag{GFT8}
 \bigotimes_{i=1}^k{}_S X(V_i)
       \ \cong\ X\left(\bigotimes_{i=1}^k{}_{\mathbb C}V_i\right),
 \qquad
 \bigotimes_i(v_i,\lambda_i)\longmapsto
              \left(\bigotimes_i v_i,\ \bigwedge_i\lambda_i\right).
\end{equation}
Parenthesizations and permutations give the same map: check them
on the displayed generators using associativity and commutativity
of meet and the corresponding ordinary tensor maps. The unit is
$X(\mathbb C)=S$. Linear maps commute with GFT8 by evaluation on
these generators. Thus this is a strong symmetric monoidal functor
on complex vector spaces with the specified lifts, not merely a
formula for selected pure tensors. Every finite sum has exactly
the join of its displayed meet labels; amplitude cancellations do
not erase any of these labels.

\subsection{Exact tensor quotients and the actual arithmetic image}

Let $W_i\subset V_i$ be arbitrary complex linear subspaces. The
ordinary tensor quotient has kernel
\begin{equation}\tag{GFT9}
 N=\sum_{i=1}^k V_1\otimes\cdots\otimes W_i\otimes\cdots\otimes V_k,
 \qquad
 \left(\bigotimes_i V_i\right)/N
                 \ \cong\ \bigotimes_i(V_i/W_i).
\end{equation}
To prove this, extend a basis of each $W_i$ to a basis of $V_i$.
The resulting algebraic tensor basis separates tensors having
at least one $W_i$ basis entry from those with only complement
entries. The first set spans exactly $N$; the second maps
bijectively to the tensor basis of the quotient spaces. The
resulting isomorphism is canonical because its formula on original
pure tensors is the quotient map, independent of those auxiliary
choices. No original tensor is replaced by the complement in its
pairings or action.

GFT8 now gives the entire kernel congruence of the lifted tensor
quotient:
\begin{equation}\tag{GFT10}
 (v,\lambda)\sim(w,\mu)
       \quad\Longleftrightarrow\quad
       \lambda=\mu\quad\hbox{and}\quad v-w\in N.
\end{equation}
Indeed two images in the constrained carrier are equal exactly
when their labels and their quotient amplitudes agree. This is
also the congruence generated by tensoring the individual
same-label quotients: GFT3's universal property forces precisely
the ordinary multilinear quotient relations on the top additive
group and leaves every lower label intact. In particular a vector
in $N$ with top label maps to $e$, not $\tau$. This congruence,
rather than the preimage of $\tau$ alone, is the relevant quotient.

For the original global adelic construction retain
\begin{equation}\tag{GFT11}
 \begin{gathered}
 C=\mathbb A_{\mathbb Q}^\times/\mathbb Q^\times,
 \quad B=\bigcap_{\beta\in\mathbb R}|\cdot|^\beta\mathcal S(C),\\
 V=\left\{x\mapsto\sum_{r\in\mathbb Q^\times}\xi(rx):
      \xi\in\mathcal S(\mathbb A_{\mathbb Q}),\
      \xi(0)=0,\ \int\xi=0\right\},\quad
 W=\overline V\subset B,\quad H=B/W.
 \end{gathered}
\end{equation}
The closure is in the original weighted Bruhat topology SWQ3;
all tensor products here remain algebraic. GFT9--10 apply to the
whole $B$, including every compact-unit level, with
$N_k=\sum_i B^{\otimes(i-1)}\otimes W\otimes
B^{\otimes(k-i)}$. Thus the entire coefficient-level tensor power
over $S$ is exactly
\begin{equation}\tag{GFT12}
 X(H)^{\otimes_S k}
     \cong X(H^{\otimes_{\mathbb C}k})
     \cong X(B^{\otimes_{\mathbb C}k})/\!\sim_{N_k}.
\end{equation}
This proves the support and amplitude map for every tensor
power of the actual quotient. It asserts no unproved topological
Kunneth theorem and does not replace $W$ by a spectral kernel.

\subsection{Null differentials before taking cohomology}

For any original cochain complex $(C^\bullet,d)$ of complex vector
spaces, retain its exact differential on each grade. Its lift obeys
\begin{equation}\tag{GFT13}
 (d^\uparrow)^2(v,\lambda)=(0,\lambda)=e(v,\lambda).
\end{equation}
For a finite tensor of such graded complexes the amplitude
differential on a homogeneous original tensor is
\begin{equation}\tag{GFT14}
 D(v_1\otimes\cdots\otimes v_k)
  =\sum_{i=1}^k(-1)^{n_1+\cdots+n_{i-1}}
       v_1\otimes\cdots\otimes d v_i\otimes\cdots\otimes v_k,
 \qquad v_i\in C_i^{n_i}.
\end{equation}
For $i<j$ the two orders of differentiating positions $i,j$
have opposite signs, since the first order increases the degree
preceding $j$ by one. Terms differentiating one position twice
vanish by the original $d^2=0$. Hence $D^2=0$ on the entire
algebraic graded tensor. GFT8 transports its lift to the actual
$S$ tensor of each homogeneous summand, preserving those signs
in amplitudes and the common meet label. In a total degree we use
the specified original synchronized carrier
$X(\bigoplus_{\sum n_i=n}\bigotimes_i C_i^{n_i})$.
It is not silently identified with a semimodule coproduct whose
summands have independent labels. The repeated sum of the one
assigned label is itself. Thus the
lifted equation is $(D^\uparrow)^2=e$, not a map constant at
$\tau$. In each total grade, original amplitude cycles with
the same-label boundary congruence give exactly
\begin{equation}\tag{GFT15}
 X(\ker D)/\!\sim_{\operatorname{im}D}
                         \cong X(\ker D/\operatorname{im}D).
\end{equation}
This follows by the same equality-of-images proof as GFT10.
It identifies the actual null complex and its cohomology map;
no nilpotent with nonzero square has been called a differential.

For precision the exact comparison with independent labels, for
any finite nonempty family of graded summands $A_j$, is
\begin{equation}\tag{GFT15a}
 \begin{gathered}
 P:\bigoplus_j X(A_j)\longrightarrow X(\bigoplus_j A_j),
       \quad ((a_j,\lambda_j))_j\longmapsto
                         ((a_j)_j,\bigvee_j\lambda_j),\\
 I:X(\bigoplus_j A_j)\longrightarrow\bigoplus_jX(A_j),
       \quad((a_j)_j,\lambda)\longmapsto((a_j,\lambda))_j,
       \qquad PI=\operatorname{id}.
 \end{gathered}
\end{equation}
All image points are in their constrained carriers: a nonzero
amplitude on either side forces the common label or at least one
input label to be top. Addition and scalar action commute with
$P,I$ by finite distributivity. The map $IP$ retains every
amplitude and replaces each input label by the displayed join.
The fibre of $P$ consists exactly of all original label tuples
with that join and the unchanged amplitudes, subject to the
original nonzero-amplitude constraints. Thus $P$ is generally
not injective. The independent-label object and its entire fibre
are retained in this comparison. GFT13--15 concern the original
common-label total complex; they do not claim that $X$ preserves
arbitrary coproducts. For an unbounded complex the synchronized
algebraic direct sum still makes sense, but GFT15a is asserted
only on each stated finite family, so no infinite join in $L$
has been assumed.

\subsection{The whole tensor pairing and its exact radical}

Let $Q$ be the actual continuous Hermitian Weil pairing on $B$,
with its full formula GTK13,24. SWQ proves $W$ is contained in
its radical, so it gives the well-defined form $Q_H$ on $H$.
Write $R=\operatorname{rad}Q_H$, retaining this subspace rather
than assuming it is zero. On the full algebraic tensor define
\begin{equation}\tag{GFT16}
 Q_H^{[k]}\left(\bigotimes_i u_i,\bigotimes_i v_i\right)
       =\prod_{i=1}^k Q_H(u_i,v_i),\qquad
 Q_H^{[k]}\left(\sum_a\bigotimes_i u_{a,i},
                    \sum_b\bigotimes_i v_{b,i}\right)
       =\sum_{a,b}\prod_i Q_H(u_{a,i},v_{b,i}).
\end{equation}
The multilinear tensor universal properties, with conjugate
linearity in the second variables, prove that these formulas are
independent of every presentation. Every factor is the entire
original global pairing, so the formula retains all cross terms
from its primes, endpoints and real place; none is replaced by
a finite window or by a completed zeta working object.

Its radical is exactly
\begin{equation}\tag{GFT17}
 \operatorname{rad}Q_H^{[k]}
      =\sum_i H^{\otimes(i-1)}\otimes R\otimes H^{\otimes(k-i)}.
\end{equation}
Inclusion of the right side is immediate from GFT16. For the
reverse inclusion pass by GFT9 to $(H/R)^{\otimes k}$, where
the original induced pairing is nondegenerate. For any finite
linearly independent list $u_1,\ldots,u_m$ in $H/R$, its pairing
functionals on the second variable are independent: a linear
relation would put a nonzero combination of the $u_i$ in the
radical. Therefore there exist $v_1,\ldots,v_m$ realizing the
coordinate dual values $Q_H(u_i,v_j)=\delta_{ij}$. This follows
also directly by considering the range in $\mathbb C^m$; a
proper range would have a nonzero annihilating functional.
Expand a nonzero algebraic tensor using independent finite lists
in each factor. Pairing with the corresponding products of
these dual vectors selects each coefficient. Some coefficient
is nonzero, proving nondegeneracy of the tensor pairing on the
quotient and hence GFT17. The argument applies to the entire
infinite-dimensional $H$, since every algebraic tensor has a
finite presentation.

The full supported form is
\begin{equation}\tag{GFT18}
 Q_{H,L}^{[k]}((u,\lambda),(v,\mu))
             =(Q_H^{[k]}(u,v),\lambda\wedge\mu).
\end{equation}
A nonzero output amplitude requires both top input labels;
linearity, conjugate linearity and finite lattice distributivity
prove the semimodule pairing identities. In particular every
top-supported tensor in the radical pairs to $e$ with a
top-supported test. GFT17 specifies precisely which amplitudes
are invisible to the whole tensor trace; their support labels
remain in GFT18.

\subsection{Exact arithmetic and support receiving maps}

There are two different scalar maps, both kept explicitly:
\begin{equation}\tag{GFT19}
 \rho:S\to\mathbb C,\ (a,\lambda)\mapsto a,
 \qquad \sigma:S\to L,\ (a,\lambda)\mapsto\lambda.
\end{equation}
They give the global base changes
\begin{equation}\tag{GFT20}
 \mathbb C\otimes_S X(V)\cong V,\quad
     c\otimes(v,\lambda)\mapsto cv;\qquad
 L\otimes_S X(V)\cong L,\quad
     \eta\otimes(v,\lambda)\mapsto\eta\wedge\lambda.
\end{equation}
For the first, the inverse is $v\mapsto1\otimes(v,1_L)$.
Every lower zero tensor vanishes because $\rho(z_\lambda)=0$;
the vector-space relations on top inputs follow from the
balanced tensor relations. These facts verify both composites
on generators. For the second, the inverse is
$\eta\mapsto\eta\otimes(0,1_L)$. Since $\sigma(e)=1_L$,
balancing gives
$\eta\otimes(v,1_L)=\eta\otimes e(v,1_L)
=\eta\otimes(0,1_L)$. A lower label reduces similarly by
balancing $z_\lambda$, giving the stated meet. This verifies
both composites without discarding any label of $L$.

Thus the arithmetic receiver of GFT12 is the entire original
$H^{\otimes k}$, while the support receiver is the entire
original $L$. The latter map has explicitly lost every amplitude,
not just the value of one function. Its fibres and a faithful
joint record are
\begin{equation}\tag{GFT21}
 X(V)\hookrightarrow V\times L,\quad(v,\lambda)\mapsto(v,\lambda),
 \qquad\operatorname{im}=
       (V\times\{1_L\})\cup(\{0\}\times L).
\end{equation}
The image condition is part of the map. The projection to $L$
has a copy of $V$ over $1_L$ and a single zero amplitude over
each other label. The projection to $V$ has a copy of the whole
$L$ over zero and one top label over each nonzero amplitude.
Consequently the joint record is injective, whereas either
projection alone need not be. For any actual support stalk map
$L\to L_J$, composition in GFT20 is exactly
$(v,\lambda)\mapsto[\lambda]_J$; it commutes with GFT8 by the
retained meet. This is the full-lattice support map, not a
replacement of the original prime spectrum by one extra point.

The obtained arithmetic tensor powers, their null differentials,
their full pairing and its radical are therefore literal objects
over the original coefficient semiring. A weight assertion must
be proved for the retained arithmetic action on these objects.
Neither the support base change GFT20 nor the algebraic tensor
isomorphism GFT8 imposes a zero-location or positive-metric
assertion. They specify the exact maps on which that further
arithmetic calculation operates.
